% Scope: Construct the phenomenological Bloch model from torque and equilibration; solve free evolution and connect reversibility, relaxation and exchange.
% Status: architecture only; no chapter prose or completed problems yet.
% Prerequisites: Chapters 1 and 2.
% Planned worked examples: Longitudinal recovery, transverse decay and an inhomogeneous ensemble.
% Planned figures: Relaxation curves and isochromat refocusing.
% Planned challenge problems: Solve a piecewise field history; distinguish static dephasing from irreversible loss.
\chapter{Classical Magnetization and the Bloch Equations}
\label{ch:03}

\section{From magnetic torque to ensemble dynamics}
\label{sec:03:from-magnetic-torque-to-ensemble-dynamics}
\subsection{Torque equation and conserved magnitude}
\subsection{Equilibrium magnetization and phenomenological relaxation}
\subsection{Assumptions of the single-pool Bloch model}

\section{Solutions in a static field}
\label{sec:03:solutions-in-a-static-field}
\subsection{Longitudinal recovery and inversion}
\subsection{Complex transverse solution and phase}
\subsection{Affine propagators and piecewise-constant fields}

\section{Inhomogeneous ensembles}
\label{sec:03:inhomogeneous-ensembles}
\subsection{Frequency distributions and free induction decay}
\subsection{T2, T2-star and reversible broadening}
\subsection{Limits of exponential effective relaxation}

\section{Driven and time-dependent dynamics}
\label{sec:03:driven-and-time-dependent-dynamics}
\subsection{Constant effective-field solution}
\subsection{Numerical integration and operator splitting}
\subsection{Accuracy, stability and conservation checks}

\section{Beyond independent stationary spins}
\label{sec:03:beyond-independent-stationary-spins}
\subsection{Bloch--Torrey transport and diffusion}
\subsection{Bloch--McConnell exchange structure}
\subsection{Where phenomenology requires microscopic input}

\section{Worked examples}
\label{sec:03:worked-examples}
% Planned calculations: Longitudinal recovery, transverse decay and an inhomogeneous ensemble.
\subsection{Derivation and quantitative calculation}
\subsection{Assumptions, limiting cases and scanner interpretation}

\section{Challenge problems}
\label{sec:03:challenge-problems}
% Problem design: Solve a piecewise field history; distinguish static dephasing from irreversible loss.
% Use challengeproblem and stable labels prob:03:<topic>.
% Complete solutions belong in Appendix E, section for this chapter.
% Use \begin{solution}{prob:03:<topic>} ... \end{solution} there.
\subsection{Derivation and quantitative design}
\subsection{Model criticism and integrated reasoning}
